\color{unidarkgreen}{\subsection{Токовая направленная защита нулевой последовательности (ТНЗНП)}}
\color{black}

\begin{enumerate}[label=\arabic{section}.\arabic{subsection}.\arabic*,  align=left, labelsep=4pt, leftmargin=0pt, itemindent=57pt]

\item
Токовая направленная защита нулевой последовательности (ТНЗНП) предназначена для селективного
отключения поврежденного участка электрической сети при коротких замыканиях на землю, реагируя на
повышение тока нулевой последовательности.

\item
В составе логики токовой направленной защиты нулевой последовательности предусмотрено шесть
ступеней от замыканий на землю.

\vspace{2mm}
\begin{figure}[H]
\centering
\includegraphics[width=1\textwidth,height=1\textheight,keepaspectratio]{\fbpath/051.0131 ТНЗНП/img/TNZNP.pdf}
\captionof{figure}{Функциональная схема ступени ТНЗНП-1}\label{pict_tnznp:TNZNP}
\end{figure}

\item
Пороги тока срабатывания ступеней ТНЗНП задаются уставками \textbf{«3I0ср ТНЗНП-1(2…6)}, а выдержки
времени срабатывания ступеней регулируются уставками \textbf{«Тср ТНЗНП-1(2…6)»}. Ввод в работу ступеней
ТНЗНП осуществляется с помощью программных переключателей \textbf{«ТНЗНП-1(2…6)»}. 

\item
Направленность ТНЗНП обеспечивается органами направления мощности нулевой последовательности
(ОНМНП) и обратной последовательности (ОНМОП).

\item
ОНМНП реагирует на величины векторов тока и напряжения нулевой последовательности, а также угол
сдвига между ними.

\item
Так как мощность нулевой последовательности всегда направлена от места несимметрии во внешнюю
сеть, то ОНМНП в АТ срабатывает при направлении мощности нулевой последовательности от АТ к шинам
(при КЗ впереди), а  ОНМНП в сеть – при обратном направлении мощности. 

\item
За положительное направление отсчета векторов токов и напряжений принято направление против
часовой стрелки. Углы максимальной чувствительности откладываются от вектора напряжения 3U0 (вектор
поляризации) и задаются уставками \textbf{«Фмч ОНМНП в АТ»} и \textbf{«Фмч ОНМНП в сеть»}. Уставки по углам
максимальной чувствительности рекомендуется задавать для ОНМНП в АТ – плюс 110º, а для ОНМНП в сеть
– плюс 290º. Ширина зоны срабатывания ± 85º от угла максимальной чувствительности.

\item
Пороги срабатывания ОНМНП в АТ и ОНМНП в сеть по току и напряжению регулируются уставками 
\textbf{«3I0ср ОНМНП в АТ»}, \textbf{«3U0ср ОНМНП а АТ»}, \textbf{«3I0ср ОНМНП в сеть»}, и \textbf{«3U0ср ОНМНП в сеть»}. Выбор
напряжения нулевой последовательности, по которому работает ОНМНП, осуществляется с помощью
программного переключателя «Выбор 3U0».

\item
Орган направления мощности обратной последовательности ОНМОП в АТ срабатывает при
направлении мощности обратной последовательности от АТ к шинам, а ОНМОП в сеть – при обратном
направлении мощности. 

\item
Углы максимальной чувствительности задаются уставками \textbf{«Фмч ОНМОП в АТ»} и \textbf{«Фмч ОНМОП
в сеть»}. Уставки по углам максимальной чувствительности рекомендуется задавать для ОНМОП в АТ – плюс
110º, а для ОНМОП в сеть – плюс 290º. Ширина зоны срабатывания ± 85º от угла максимальной
чувствительности.

\item
Пороги срабатывания ОНМОП в АТ и ОНМОП в сеть по току и напряжению регулируются уставками 
\textbf{«3I0ср ОНМОП в АТ»}, \textbf{«3U0ср ОНМОП в АТ»}, \textbf{«3I0ср ОНМОП в сеть»}, и \textbf{«3U0ср ОНМОП в сеть»}.

\small
\setlength{\extrarowheight}{0.05cm} % высота ячеек в таблице
\begin{longtable}
{|>{\centering\arraybackslash}m{15.2cm}|>{\centering\arraybackslash}m{1.7cm}|}
\caption{Параметры ОНМНП(ОНМОП)\hfill\vspace{-0.5\baselineskip}}\label{desc_tnznp:tbl1}\\ 
\hhline{|-|-|}
\rowcolor{gray!30}
Наименование параметра & Значение  \\ 
\hhline{|-|-|}
\endfirsthead

\caption*{\hspace{3pt}\emph{Продолжение таблицы \ref{desc_tnznp:tbl1}\hfill\vspace{-0.5\baselineskip}}} \\ % сделано по ГОСТ 2.105 п.6.8.7
\hhline{|-|-|}
\rowcolor{gray!30}
Наименование параметра & Значение  \\ 
\hhline{|-|-|}
\endhead
\hhline{|-|-|}

\endfoot
\endlastfoot

\raggedright Время срабатывания ОНМ при утроенных по отношению к порогам срабатывания значениях тока и напряжения, мс, не более & \centering\arraybackslash 35  \\  \hline
\raggedright Время возврата ОНМ при одновременном сбросе входных величин тока и напряжения от номинальных значений до нуля, мс, не более  & \centering\arraybackslash 25  \\  \hline
\raggedright Средняя основная погрешность порогов срабатывания ОНМ по току и напряжению, \% уставки, не более   & \centering\arraybackslash $\pm$5  \\  \hline
\raggedright Средняя основная абсолютная погрешность ОНМ по углу максимальной чувствительности, эл. град., не более & \centering\arraybackslash $\pm$5 \\  \hline
\raggedright Дополнительная погрешность параметров срабатывания от изменения температуры окружающего воздуха в рабочем диапазоне, \% от среднего значения параметров, измеренных при температуре $(20\pm5)^\circ$C, не более & \centering\arraybackslash $\pm$5 \\  \hline
\raggedright Минимальная угловая ширина зоны срабатывания ОНМ, эл. град, не менее & \centering\arraybackslash 165 \\ \hline
\raggedright Коэффициент возврата ОНМ НП по напряжению и току, не менее  & \centering\arraybackslash 0,9\\  \hline

\end{longtable}
\setlength{\extrarowheight}{0.0cm} % возврат параметра высота ячеек в таблице
\normalsize

\item
На рисунке \ref{pict_tnznp:char_triggering_ONMNP} представлены характеристики срабатывания ОНМНП в АТ и ОНМНП в сеть. Аналогичные
характеристики для ОНМОП в АТ и ОНМОП в сеть.

\vspace{2mm}
\begin{figure}[H]
\centering
\includegraphics[width=1\textwidth,height=1\textheight,keepaspectratio]{\fbpath/051.0131 ТНЗНП/img/char_triggering_ONMNP.pdf}
\captionof{figure}{Характеристики срабатывания ОНМНП}\label{pict_tnznp:char_triggering_ONMNP}
\end{figure}

\item
Для выбора режима направленности по мощности нулевой последовательности для каждой ступени
ТНЗНП предусмотрен программный переключатель \textbf{«Направл ТНЗНП-1(2…6)»}, имеющий три положения:

\begin{enumerate}[label=\arabic*)]
    \item \textit{«Ненапр»} -- ступень выполнена ненаправленной;
    \item \textit{«В АТ»} -- ступень выполнена прямонаправленной в сторону АТ;
    \item \textit{«В сеть»} -- ступень выполнена обратнонаправленной в сторону сети.
\end{enumerate}

\item
Кроме того, можно установить различные алгоритмы контроля направленности с помощью программных
переключателей \textbf{«Контр направл ТНЗНП-1(2…6)»}, имеющих три положения:

\begin{enumerate}[label=\arabic*)]
    \item \textit{«Разр»} --  прямая направленность обеспечивается срабатыванием ОНМНП в АТ, обратная -- срабатыванием ОНМНП в сеть;
    \item \textit{«Блок»} -- прямая направленность обеспечивается отсутствием срабатывания ОНМНП в сеть, обратная -- отсутствием срабатывания ОНМНП в АТ;
    \item \textit{«Разр или блок»} -- прямая направленность обеспечивается срабатыванием ОНМНП в АТ или отсутствием срабатывания ОНМНП в сеть, обратная -- срабатыванием ОНМНП в сеть или отсутствием срабатывания ОНМНП в АТ.
\end{enumerate}

\item
Положение переключателя \textit{«Разр. или блок.»} рекомендуется выставлять для ступеней ТНЗНП,
выполняющих функцию дальнего резервирования, так как при удаленных КЗ напряжение $3U_0$ может оказаться
недостаточным для срабатывания ОНМ из-за его уменьшения по мере удаления от точки КЗ.

\item
Уставки для выбора режима направленности по мощности обратной последовательности аналогичные.

\item
Тип ОНМ, используемый для определения направленности ступени ТНЗНП, выбирается программным
переключателем \textit{«Тип ОНМ ТНЗНП-1(2…6)»}, имеющим четыре положения:

\begin{enumerate}[label=\arabic*)]
    \item \textit{«НП»} -- для определения направленности ступени учитывается только ОНМНП;
    \item \textit{«ОП»} -- для определения направленности ступени учитывается только ОНМОП;
    \item \textit{«НП или ОП»} -- для определения направленности ступени учитывается один из органов, ОНМНП или ОНМОП;
    \item \textit{«НП и ОП»} -- для определения направленности ступени одновременно учитываются оба органа, ОНМНП и ОНМОП.
\end{enumerate}

\item
Имеется возможность задать различные режимы работы ступеней ТНЗНП при возникновении
неисправности в цепях напряжения. Вариант действия логики направленных ступеней ТНЗНП при
срабатывании БНН задается программным переключателями \textbf{«Режим БНН ТНЗНП-1(2…6)»}, имеющими три
положения:

\begin{enumerate}[label=\arabic*)]
    \item \textit{«Выведено»} -- режим работы ступени не зависит от БНН;
    \item \textit{«ОНМ»} -- ступень переводится в ненаправленный режим работы;
    \item \textit{«ТНЗНП»} -- при срабатывании БНН ступень выводится из работы. Однако, если ступень задана ненаправленной, то она останется в работе независимо от срабатывания БНН. Ступень также не будет выводиться из работы, если направленность автоматически вывелась при срабатывании ТНЗНП.
\end{enumerate}

\item
Предусмотрен автоматический вывод направленности всех ступеней ТНЗНП при срабатывании ТНЗНП,
это обеспечивает устойчивое срабатывание ТНЗНП при неполнофазном отключении, что необходимо для
надежного действия УРОВ. Данный режим вводится с помощью программного переключателя \textbf{«Вывод направл при ср ТНЗНП»}.

\item
Для обеспечения несрабатывания ТНЗНП от бросков намагничивающего тока предусмотрена
блокировка ступеней ТНЗНП по содержанию второй гармоники в токе нулевой последовательности.
Блокировка основана на контроле отношения действующего значения второй гармоники тока нулевой
последовательности к основной гармонике. Ступени ТНЗНП блокируется при превышении данного
соотношения пороговой величины, заданной уставкой \textbf{«3I02г/3I01г ТНЗНП»}. Наличие функции блокировки от
бросков намагничивающего тока задается для каждой ступени отдельно с помощью программных
переключателей \textbf{«Режим БНТ ТНЗНП-1(2…6)»}. Для предотвращения излишней блокировки при больших токах
КЗ, сопровождающихся насыщением ТТ, блокировка выводится из действия при превышении тока 3I0
основной гармоники выше уставки \textbf{«3I0блок ОВ БНТ ТНЗНП»}.

\item
Оперативный ввод/вывод чувствительных ступеней ТНЗНП производится с помощью команды \textit{«Вывод
ЧС ТНЗНП»}, а выбор выводимых ступеней ТНЗНП определяется уставками \textbf{«Режим БЧС ТНЗНП-1(2…6)»}.

\item
Оперативный ввод/вывод ТНЗНП производится с помощью БУ (блоков управления функциями) \textit{«Вывод ТНЗНП»}. 

\item
Пуски степеней ТНЗНП направленные в сеть, действуют в логику предварительного деления своей
тороны АТ. С помощью выдержек времени в логике деления \textbf{«Тср ТНЗНП-1(2…6) на СВ»}, \textbf{«ТНЗНП-1(2…6) на
В сеч»} и \textbf{«Тср ТНЗНП-1(2…6) на В»}, первым действием формируются сигналы на разделение секции
(отключение СВ) или деление по сечению (отключение В сечения), вторым действием формируются сигналы
на отключение выключателей своей стороны с пуском УРОВ и запретом АПВ. Пуски ступеней ТНЗНП
направленные в АТ, могут действовать в логику предварительного деления резервных защит смежной стороны
(смежный комплект РЗАТ), через выдержки времени в логике деления \textbf{«Тср ТНЗНП-1(2…6) на пуск ЛД см
ст»}. Действие срабатывания ступеней ТНЗНП предусматривает отключение АТ со всех сторон с запретом
АПВ и пуском УРОВ.

\item
В составе логики токовой направленной защиты нулевой последовательности предусмотрено четыре
режима оперативного ускорения ТНЗНП.

\vspace{2mm}
\begin{figure}[H]
\centering
\includegraphics[width=0.9\textwidth,height=0.9\textheight,keepaspectratio]{\fbpath/051.0131 ТНЗНП/img/OA_TNZNP.pdf}
\captionof{figure}{Функциональная схема ОУ ТНЗНП}\label{pict_tnznp:OA_TNZNP}
\end{figure}

\item
Оперативное ускорение ТНЗНП в сеть без выдержки времени. Предназначено для оперативного
ускорения ступени ТНЗНП  при выполнении переключений в сети с действием на отключение всего АТ (с пуском
УРОВ и запретом АПВ) с одной выдержкой времени, которая регулируется уставкой \textbf{«Тср ОУ ТНЗНП»}, выбор
ускоряемой ступени определяется программным переключателем \textbf{«ОУ ТНЗНП»}. Оперативный ввод ОУ ТНЗНП
в сеть без выдержки времени производится с помощью команды управления \textit{«Ввод ОУ ТНЗНП в сеть без вв»}. 

\item
Оперативное ускорение ТНЗНП в сеть с первой выдержкой времени. Предназначено для оперативного
ускорения ступени ТНЗНП при выводе ДЗО с действием на отключение всего АТ (с пуском УРОВ и запретом
АПВ) с одной выдержкой времени, которая регулируется уставкой \textbf{«Тср ОУ1 ТНЗНП»}, выбор ускоряемой
ступени определяется программным переключателем \textbf{«ОУ1 ТНЗНП»}. Оперативный ввод ОУ ТНЗНП в сеть с
первой выдержкой времени производится с помощью команды управления \textit{«Ввод ОУ ТНЗНП в сеть с вв1»}.

\item
Оперативное ускорение ТНЗНП в сеть со второй выдержкой времени. Предназначено для
оперативного ускорения ступени ТНЗНП при выводе ДЗШ с действием пуска ОУ2 ТНЗНП в логику
предварительного деления своей стороны. Действие срабатывания ОУ2 ТНЗНП предусматривает отключение
всего АТ (с пуском УРОВ и запретом АПВ). Время срабатывания ОУ2 ТНЗНП регулируется уставкой \textbf{«Тср ОУ2
ТНЗНП»}, а выбор ускоряемой ступени определяется программным переключателем \textbf{«ОУ2 ТНЗНП»}.
Оперативный ввод ОУ ТНЗНП в сеть со второй выдержкой времени производится с помощью команды
управления \textit{«Ввод ОУ ТНЗНП в сеть с вв2»}.

\item
Оперативное ускорение ТНЗНП в АТ. Предназначено для оперативного ускорения ступени ТНЗНП при
выводе основных защит АТ с действием на отключение всего АТ (с пуском УРОВ и запретом АПВ) с одной
выдержкой времени, которая регулируется уставкой \textbf{«Тср ОУ ТНЗНП в АТ»}, выбор ускоряемой ступени
определяется программным переключателем \textbf{«ОУ ТНЗНП в АТ»}. Оперативный ввод ОУ ТНЗНП в АТ
производится с помощью команды управления \textit{«Ввод ОУ ТНЗНП в АТ»}.

\item
В составе логики токовой направленной защиты нулевой последовательности предусмотрен пуск
автоматического ускорения ступеней ТНЗНП.

\vspace{2mm}
\begin{figure}[H]
\centering
\includegraphics[width=0.9\textwidth,height=0.9\textheight,keepaspectratio]{\fbpath/051.0131 ТНЗНП/img/AA_TNZNP.pdf}
\captionof{figure}{Функциональная схема пуска АУ ТНЗНП}\label{pict_tnznp:AA_TNZNP}
\end{figure}

\item
С помощью программного переключателя \textbf{«Пуск АУ ТНЗНП в АТ»} имеется возможность выбрать от
какой ступени ТНЗНП будет пускаться автоматическое ускорение при опробовании АТ.

\item
С помощью программного переключателя \textbf{«Пуск АУ ТНЗНП в сеть»} имеется возможность выбрать от
какой ступени ТНЗНП будет пускаться автоматическое ускорение при опробовании ошиновки АТ или системы шин.

\item
Дополнительно предусмотрена возможность пуска ускорения от ненаправленного ИО ступени ТНЗНП c
контролем броска намагничивающего тока, который задается программным переключателем \textbf{«Режим БНТ АУ
ТНЗНП ненапрв»}, а так же помощью программного переключателя \textbf{«Пуск АУ ТНЗНП ненапрв»} имеется
возможность выбрать от какого ненаправленного ИО ступени будет пускаться автоматическое ускорение при
включении выключателя. 

\item
В составе логики токовой направленной защиты нулевой последовательности предусмотрен режим
ускорения от защит смежной стороны.

\vspace{2mm}
\begin{figure}[H]
\centering
\includegraphics[width=0.9\textwidth,height=0.9\textheight,keepaspectratio]{\fbpath/051.0131 ТНЗНП/img/TA_TNZNP.pdf}
\captionof{figure}{Функциональная схема ускорения ТНЗНП от защит смежной стороны}\label{pict_tnznp:TA_TNZNP}
\end{figure}

\item
Предусматривается возможность встречного ускорения ТНЗНП от сигнала пуска ступени ТНЗНП \textit{«Уск.
от ТНЗНП см. ст»}, смежного комплекта резервных защит автотрансформатора (РЗАТ). Так же, с помощью
программного переключателя \textbf{«Уск ТНЗНП от РПО см ст»} предусматривается возможность разрешения
встречного ускорения ТНЗНП при наличии сигнала отключенного положения выключателей смежной стороны
\textit{«См. ст. отключена»}. Данный вид ускорения позволяет использовать ТНЗНП как защиту абсолютной
селективности, так как пуски защит направленных в АТ на комплектах РЗАТ ВН и СН указывают на КЗ внутри
АТ. С помощью программного переключателя \textbf{«Уск см ТНЗНП»} имеется возможность выбрать от пуска какой
ступени ТНЗНП будет контролироваться ускорение от защит смежной стороны, а также формироваться сигнал
пуска ускорения в защиты смежной стороны \textit{«Уск. ТНЗНП на см. ст»}. Действие защиты предусматривает
отключение АТ со всех сторон с пуском УРОВ и запретом АПВ, минуя этапы деления. Выдержка времени
срабатывания ускорения ТНЗНП от защит смежной стороны регулируется уставкой \textbf{«Тср уск ТНЗНП от см ст»}.

\small
\begin{longtable}{|>{\centering\arraybackslash}m{5.3cm}|>{\centering\arraybackslash}m{3.3cm}|>{\centering\arraybackslash}m{4.2cm}|>{\centering\arraybackslash}m{1.8cm}|>{\centering\arraybackslash}m{1cm}|}
\caption{Перечень уставочных параметров ТНЗНП\hfill\vspace{-0.5\baselineskip}}\label{tnznp:tbl1}\\ 
\hline
\rowcolor{gray!30}
Параметр & Условное обозначение на схеме & Значение/ Диапазон & Единица измерения & Шаг \\ 
\hline
\endfirsthead
\caption*{\hspace{3pt}\emph{Продолжение таблицы \ref{tnznp:tbl1}\hfill\vspace{-0.5\baselineskip}}} \\ % сделано по ГОСТ 2.105 п.6.8.7
\hline
\rowcolor{gray!30}
Параметр & Условное обозначение на схеме & Значение/ Диапазон & Единица измерения & Шаг \\ 
%\hline
\endhead

%===m>ТНЗНП|-
%===m

\end{longtable}
\normalsize

\end{enumerate}